\documentclass[11pt,a4paper]{article}
\usepackage{fontspec,amsmath,amssymb,amsthm,microtype}
\setmainfont{lmroman10-regular.otf}[BoldFont=lmroman10-bold.otf,ItalicFont=lmroman10-italic.otf,BoldItalicFont=lmroman10-bolditalic.otf]
\usepackage{xcolor}
\usepackage[most]{tcolorbox}
\usepackage{mathrsfs}
\providecommand{\lt}{<}
\providecommand{\gt}{>}
\setmonofont{DejaVuSansMono.ttf}[BoldFont=DejaVuSansMono-Bold.ttf,ItalicFont=DejaVuSansMono-Oblique.ttf,BoldItalicFont=DejaVuSansMono-BoldOblique.ttf,Scale=0.86]
\newfontfamily\leanfallback{STIXTwoMath-Regular.otf}[Scale=0.86]
\definecolor{leanrule}{gray}{0.62}
\definecolor{leanback}{gray}{0.962}
\definecolor{leantitle}{gray}{0.22}
\definecolor{leanmuted}{gray}{0.38}
\newtcolorbox{leanbox}[1]{enhanced,breakable,colback=leanback,colframe=leanrule,
  boxrule=0.4pt,arc=0.6mm,left=2.4mm,right=2.4mm,top=1.8mm,bottom=1.8mm,
  before skip=0.8em,after skip=1em,fonttitle=\small\bfseries,coltitle=leantitle,
  colbacktitle=leanback,titlerule=0.3pt,title={#1}}
\newlength{\leancharwidth}
\newenvironment{leancode}{\par\smallskip\begingroup\ttfamily\footnotesize\raggedright
  \settowidth{\leancharwidth}{M}\setlength{\parskip}{0pt}\setlength{\parindent}{0pt}
  \frenchspacing}{\par\endgroup\smallskip}
\newcommand{\leanline}[2]{\hangindent=\dimexpr#1\leancharwidth+4\leancharwidth\relax
  \hangafter=1\noindent\hspace*{#1\leancharwidth}#2\par}
\newcommand{\leansource}[1]{\par{\footnotesize\color{leanmuted}#1}\par}
\newcommand{\leanlabel}[1]{\par\smallskip\noindent{\small\bfseries #1}\par}
\usepackage[margin=25mm]{geometry}
\usepackage[unicode,hidelinks]{hyperref}
\hypersetup{pdfauthor={LangevinLean Project}}
\setlength{\emergencystretch}{3em}
\setlength{\parskip}{0.4em}
\setlength{\parindent}{0pt}
\allowdisplaybreaks[2]
\raggedbottom
\title{\texorpdfstring{Conditional expectation of a genuine integrated defect}{Conditional expectation of a genuine integrated defect}}
\author{LangevinLean Project}
\date{}
\begin{document}\maketitle
Draft mathematical exposition; these notes may contain mistakes.


The actual Hilbert conditional-expectation projection commutes with Bochner time integration. If each increment has conditional mean equal to an evaluated source, the time-integrated defect has conditional mean equal to the evaluated integrated source.


\section{\texorpdfstring{Commuting the actual bounded projections and time integral}{Commuting the actual bounded projections and time integral}}


\textbf{Claim.} If \(\mathbb E[I(t)\mid\mathcal F]=U s(t)\) for \(0\le t\le h\), where \(U\) is an evaluation isometry and the increment curve is integrable, then

\[\mathbb E\left[\int_0^h I(t)\,dt\,\middle|\,\mathcal F\right]=U\left(\int_0^h s(t)\,dt\right).\]

\textbf{Proof.} Express conditional expectation as the orthogonal projection onto the closed past-measurable \(L^2\) subspace. It is bounded linear, so it commutes with the justified integral. Substitute the pointwise conditional-mean identities inside the integral. Evaluation is also bounded linear and commutes with the same integral, yielding the displayed result.

\begin{leanbox}{Lean: \texttt{conditionalDefect\_\allowbreak{}eq\_\allowbreak{}evaluation\_\allowbreak{}integral}}
\leansource{\texttt{conditionalDefect\_\allowbreak{}eq\_\allowbreak{}evaluation\_\allowbreak{}integral} \quad WassersteinMixing/\allowbreak{}ConditionalDefect.lean, lines 24–39}
\begin{leancode}
\leanline{0}{theorem\ conditionalDefect\_eq\_evaluation\_integral}
\leanline{4}{(eval\ :\ R\ →ₗᵢ[ℝ]\ Lp\ E\ 2\ Q)\ (increment\ :\ ℝ\ →\ Lp\ E\ 2\ Q)\ (source\ :\ ℝ\ →\ R)}
\leanline{4}{(h\ :\ ℝ)\ (hh\ :\ 0\ ≤\ h)\ (hi\ :\ IntervalIntegrable\ increment\ volume\ 0\ h)}
\leanline{4}{(hmean\ :\ ∀\ t\ ∈\ Set.Icc\ 0\ h,}
\leanline{6}{(condExpL2\ E\ ℝ\ hm\ (increment\ t)\ :\ Lp\ E\ 2\ Q)\ =\ eval\ (source\ t))\ :}
\leanline{4}{(condExpL2\ E\ ℝ\ hm\ (stationaryLocalDefect\ increment\ h)\ :\ Lp\ E\ 2\ Q)\ =}
\leanline{6}{eval\ (∫\ t\ in\ 0..h,\ source\ t)}
\end{leancode}
\leanlabel{How the Lean proof follows the argument}
\texttt{conditional\allowbreak{}Defect\_\allowbreak{}eq\_\allowbreak{}evaluation\_\allowbreak{}integral} uses the actual \texttt{lp\allowbreak{}Meas} orthogonal projection and its interval-integral commutation theorem, then the evaluation isometry integral identity. Pointwise substitution is justified on the full closed time interval.
\end{leanbox}

\end{document}
